\chapter{What MeerK40t does, and what you need}
\label{ch:intro}

Before you touch a menu, it is worth being clear about what this program is for
and what the words mean. This chapter is the map; the rest of the manual is the
walk.

\section{The one-paragraph version}

MeerK40t is a free, open-source program that drives laser cutters and engravers.
You give it artwork --- a drawing, a photograph, a shape you drew in the program
itself --- tell it how that artwork should be burned, and it converts both into
the precise stream of movements and laser pulses your machine understands. It
runs on Windows, macOS, Linux and the Raspberry~Pi, supports several popular
controller families, and is extensible through plugins. It is free software under
the MIT licence.

\begin{behindbox}[Behind the name and the spelling]
``MeerK40t'' is a pun on the K40, the inexpensive Chinese CO\textsubscript{2}
laser cutter that the project originally targeted, and on ``meerkat''. ``K40''
with ``t'' appended gives ``K40t'', read as the sound in ``meerkat''. You will
see the name written both \emph{MeerK40t} (the program) and \emph{Meerk40t}
(variant capitalisation in older text, screenshots and file names). It is the
same thing.
\end{behindbox}

\section{The three things the program does}

Everything in MeerK40t belongs to one of three activities, and every chapter of
this manual is about one of them.

\begin{enumerate}[leftmargin=1.6em]
  \item \textbf{Hold artwork} --- shapes, paths, text and images live in a
        document called the \emph{element tree} (Chapter~\ref{ch:tree} covered
        later). This is where you import, arrange, edit and position.
  \item \textbf{Describe the burn} --- each group of artwork is attached to an
        \emph{operation}: cut it, engrave it, raster it, dot it. Operations carry
        the machine settings: speed, power, passes, resolution and so on.
  \item \textbf{Run the job} --- MeerK40t turns operations into
        \emph{cutcode}, optimises the order of the moves, wraps it in a
        \emph{job}, hands it to a spooler and streams it to the controller.
\end{enumerate}

\begin{figure}[htbp]
\centering
\begin{tikzpicture}[
  node distance=4mm,
  box/.style={draw=mkblue!70,fill=mkblue!6,rounded corners=1pt,
              text width=1.95cm,align=center,font=\sffamily\scriptsize,
              minimum height=1.0cm,inner sep=3pt},
  stage/.style={draw=mkgrey!60,fill=black!4,rounded corners=1pt,
              text width=1.4cm,align=center,font=\sffamily\scriptsize,
              minimum height=0.85cm,inner sep=2pt},
  arrow/.style={-{Stealth[length=2mm]},draw=mkdark,thick}]
\node[box] (art) {Artwork\\\scriptsize(element tree)};
\node[box,right=7mm of art] (ops) {Operations\\\scriptsize(cut / engrave /\\raster / dots)};
\node[box,right=7mm of ops] (code) {Cutcode\\\scriptsize(lines, raster\\ strokes, moves)};
\node[box,right=7mm of code] (job) {Job\\\scriptsize(spooler)};
\node[box,right=7mm of job] (dev) {Device\\\scriptsize(driver + USB / net)};
\draw[arrow] (art) -- (ops);
\draw[arrow] (ops) -- (code);
\draw[arrow] (code) -- (job);
\draw[arrow] (job) -- (dev);

\node[stage,below=8mm of code,xshift=-1.5cm] (copy) {copy};
\node[stage,right=1mm of copy] (pre) {preprocess};
\node[stage,right=1mm of pre] (val) {validate};
\node[stage,right=1mm of val] (blob) {blob};
\node[stage,right=1mm of blob] (op1) {preopt};
\node[stage,right=1mm of op1] (op2) {optimize};
\node[stage,right=1mm of op2] (spool) {spool};
\node[font=\sffamily\scriptsize,below=3mm of val,color=mkgrey,
  text width=9cm,align=center]
  {the six planning stages, plus spool, named in the Console output whenever a job is built};
\draw[draw=mkgrey!70,dashed] ($(copy.north west)+(-2mm,2.5mm)$) rectangle ($(spool.south east)+(2mm,-1.5mm)$);
\end{tikzpicture}
\caption{The path from artwork to laser. The upper row is the journey you will
follow all manual long; the lower row is the planning pipeline described in
Chapter~\ref{ch:order} and printed in the console whenever a job is built.}
\label{fig:pipeline}
\end{figure}

\section{Laser vocabulary}

You can use MeerK40t without this section, but you will misread the screens
without it. Learn these ten terms.

\begin{description}[leftmargin=1.4em,style=nextline]
  \item[Vector] A shape described mathematically: a line, an arc, a bezier
        curve. Vectors scale without pixellating. SVG and DXF files are vector
        files.
  \item[Raster] A grid of pixels, as in a photograph or a PNG. Engraving a
        raster means burning it line by line, switching the beam on and off (or
        varying its power) as the head sweeps across.
  \item[Cut / Score] A vector burned at high power and low speed, through the
        material. A \emph{score} or \emph{mark} is the same path burned lightly
        so it is visible but not severed --- a fold line, for example. In
        MeerK40t both are done with a \mkseries{Cut} or \mkseries{Engrave}
        operation, differing only in settings.
  \item[Engrave] A vector burned at engraving settings: the beam traces the
        path, removing a thin line of material. Dark, thin, precise --- but slow
        over a large area. Filling an area by engraving a dense set of lines is
        called \emph{hatching}.
  \item[Speed and power] The two dials that decide everything. Speed is how fast
        the head moves (millimetres per second, or per minute on some machines);
        power is how hard the beam burns. Thicker material wants slower and/or
        more power.
  \item[Passes] How many times the beam repeats the same path. Two passes at
        moderate power cut cleaner than one at maximum, especially on plywood
        and acrylic.
  \item[PPI / frequency] Pulses Per Inch --- how many laser pulses per inch of
        travel. Older controllers expose pulse rate; modern ones expose
        frequency in kHz. It changes how a cut edge looks more than whether it
        cuts.
  \item[Kerf] The width of material the beam destroys. A 0.15\unit{mm} kerf
        means a part designed 20\unit{mm} wide comes out 19.85\unit{mm} wide
        unless you compensate. MeerK40t can offset paths to compensate.
  \item[DPI / raster step] Dots per inch --- how many scan lines per inch in a
        raster engrave. More DPI is finer and much slower. MeerK40t usually
        works in \emph{raster step} internally (the distance between lines) and
        shows DPI beside it.
  \item[Origin / home / jog] \emph{Home} is a fixed corner the machine can find
        itself (usually by limit switches). \emph{Origin} is the zero point the
        job is measured from. \emph{Jogging} is moving the head by hand without
        cutting.
\end{description}

\begin{notebox}[Focus]
Focus --- the height of the lens above the material --- is not controlled by
MeerK40t on most hobby machines: you set it with a gauge or a motorised bed.
Everything in this program assumes the focus is correct. Blurry cuts with
otherwise sensible settings almost always mean focus, not software.
\end{notebox}

\section{Which machines are supported}

MeerK40t drives machines through dedicated \emph{drivers}. Find your controller
in the table below; if you do not know which one you have, Chapter~\ref{ch:connect}
helps you identify it and the manufacturer's model name is usually enough.

\begin{center}\small
\begin{tabular}{@{}L{2.6cm}L{5.4cm}L{5.9cm}@{}}
\toprule
\textbf{Driver} & \textbf{Typical machines} & \textbf{Notes} \\
\midrule
Lihuiyu & K40 and other M2-Nano / M3-Nano boards & USB, no controller-side
settings, raster and vector both supported \\
GRBL & Diode engravers: Ortur, Atomstack, Longer, Creality, Sculpfun, generic
G-code machines & The most predictable driver; serial or USB, also network
(Chapter~\ref{ch:remote}) \\
Ruida & Black-and-yellow CO\textsubscript{2} machines with an RDC6442/RDC6445
panel & Network-first, and also able to \emph{emulate} Ruida for other software \\
Moshi & Older Moshiboard machines & USB \\
Newly & NewlyDraw System 8.1 machines & USB, PWM frequency exposed \\
Balor (galvo) & Fibre and UV galvo markers with Ezcad2-compatible JCZ controllers
& Field-based marking, not gantry motion \\
\bottomrule
\end{tabular}
\end{center}

\begin{center}\small
\begin{tabular}{@{}L{2.6cm}L{3.2cm}L{3.9cm}L{4.2cm}@{}}
\toprule
\textbf{Driver} & \textbf{Use status} & \textbf{Development} & \textbf{Known issues} \\
\midrule
Lihuiyu & usable & under revision & reported raster Y-registration issues \\
GRBL & good & stable & none known \\
Moshi & good & stable & none known \\
Newly & good & stable & none known \\
Ruida & good & under revision & fit for testing rather than production \\
Balor & good & stable & none known \\
\bottomrule
\end{tabular}
\end{center}

\noindent That second table is worth reading carefully if your work matters: this
manual describes features that are complete and supported, but a driver marked
``under revision'' is exactly what it says. Test any new machine profile on
scrap before running something expensive.

\section{What MeerK40t is not}

Setting expectations saves frustration:

\begin{itemize}
  \item It is \textbf{not a CAD or illustration program}. It can draw shapes,
        edit paths and lay out text, and many people do whole projects in it ---
        but serious artwork is usually made in Inkscape, Illustrator, Affinity
        Designer or Fusion~360 and imported (Chapter~\ref{ch:files}).
  \item It is \textbf{not a photo editor}. It rasterises photographs beautifully,
        but cropping, retouching and levels belong in an image editor first.
  \item It does \textbf{not replace your machine's manual}. Homing behaviour,
        limit switches, interlocks, water chillers and lens maintenance are
        machine matters.
\end{itemize}

\section{What you need before Chapter~\ref{ch:firstcut}}

\begin{itemize}
  \item A supported laser, powered, with its interlocks intact.
  \item The cable your machine uses (usually USB; Ruida machines are usually
        network).
  \item MeerK40t installed and starting (Chapter~\ref{ch:install}).
  \item A few offcuts of the material you intend to use. Your first job should
        burn scrap, not your final workpiece.
  \item A way to secure the workpiece: tape, magnets, or the machine's own
        clamps. Nothing moves a laser cut like a workpiece that shifted
        mid-job.
  \item Air assist, if your machine has it, and a focused lens.
\end{itemize}

\begin{tryit}{Inventory your machine before you start}
Write down, in one place, the answers to these questions. You will need every
one of them when you configure the device in Chapter~\ref{ch:connect}:

\begin{enumerate}[label=\arabic*.,leftmargin=1.4em]
  \item What is the working area of the bed, in millimetres? (Measure it; the
        specification is often optimistic.)
  \item Where is the machine's origin --- front-left, front-right, or centre?
  \item What controller is inside it --- Lihuiyu, GRBL, Ruida, Moshi, Newly,
        or a galvo head?
  \item What is the maximum power of the tube or diode, in watts?
  \item Does it home itself on power-up, and in which corner?
  \item Does it have air assist, a rotary attachment, or a camera?
\end{enumerate}
\end{tryit}
